\chapter*{Canonical Notation Register}
\markboth{Canonical Notation Register}{Canonical Notation Register}

\begingroup

\small
\setlength{\LTleft}{0pt}
\setlength{\LTright}{0pt}
\setlength{\tabcolsep}{1.5pt}
\renewcommand{\arraystretch}{1.10}

\newcommand{\CanonicalRef}[1]{%
  \hyperref[#1]{\ref*{#1}, p.~\pageref*{#1}}%
}

\newcommand{\CanonicalEntry}[6]{%
  #1 & #2 & #3 & #4 & #5 & \CanonicalRef{#6}\\
}

\newcommand{\CanonicalLaw}[3]{%
  #1 & #2 & \CanonicalRef{#3}\\
}

\newcommand{\CanonicalOverload}[3]{%
  #1 & #2 & #3\\
}

\newcommand{\CanonicalTerm}[4]{%
  #1 & #2 & #3 & \CanonicalRef{#4}\\
}

\newcommand{\CanonicalTermOverload}[3]{%
  #1 & #2 & #3\\
}

\section*{Register policy}

This file is the author-facing source of truth for persistent notation and technical
terminology introduced or fixed by the spectral algebraic geometry chapter.  It records
symbols together with their mathematical type, canonical meaning, scope, publication
status, and source anchor, and it records terms whose technical meaning must remain stable
in later chapters.  It is not a prose glossary: ordinary locally bound algebra objects,
modules, points, covering indices, and diagram variables are omitted unless the chapter
promotes them to persistent notation or a later interface depends on their exact use.

All categories are $\infty$-categories unless explicitly stated otherwise.  Mapping spaces
are written $\Map_{\mathcal C}(X,Y)$ and mapping spectra in stable categories are written
$\map_{\mathcal C}(X,Y)$.  Coherent higher structure produced by functors, adjunctions,
coCartesian transport, operadic constructions, Kan extensions, limits, and descent
totalizations is the canonical structure supplied by those constructions; it is not
supplemented by independently chosen coherence data.

The basic geometric variance is fixed once and for all by
\[
  \Aff_{\Sp}=(\CAlg(\Sp)^{\mathrm{cn}})^{\op}.
\]
Thus a geometric map $f:Y\to X$ reindexes coefficients contravariantly by $f^*$.
On affine schemes this is extension of scalars, with restriction of scalars as right
adjoint.  The same symbols $f^*$ and $f_*$ are retained globally because the affine and
global constructions are identified by the right-Kan-extension and descent results.

Classical truncation is likewise transported compatibly through the chapter.  On affines
$\Spec(A)$ it is $\Spec\pi_0(A)$; on prestacks it is restriction $\delta^*$ to discrete affine
spectral schemes; on spectral schemes it lands in ordinary schemes.  The notation
$(-)_{\mathrm{cl}}$ and $X_{\mathrm{cl}}$ is deliberately reused for these compatible
levels.  For an ordinary commutative ring $R$, $\Spec(R)$ denotes the ordinary affine scheme,
while $\Spec(HR)$ denotes its discrete affine spectral enhancement.

The chapter distinguishes stable exactness from geometric flatness.  Extension of scalars
is exact as a functor of stable $\infty$-categories for every algebra map, while flatness
refers to $t$-exactness for the relevant Postnikov structures.  Similarly, ``finite cell''
denotes an explicit finite-colimit construction, whereas ``locally of finite presentation''
denotes compactness in the appropriate undercategory.

The notations $\mathbb L_X(\mathcal E)$ and $\mathbf V_X(\mathcal E)$ are intentionally
distinct.  The former is the relative linear space associated directly to a connective
quasi-coherent module; the latter is the section total space of a finite locally free
module and is defined by dualizing first.

References in this register use existing stable labels from the spectral algebraic geometry
chapter.  No cross-reference label is invented only inside a notation file.

In the \textbf{Public} column, ``yes'' means that the canonical symbol is eligible for
projection to the publication-facing List of Notation.  The public list is a lean navigation
index and may merge tightly coupled entries or append a registered defining equation, but
every head notation item displayed there is backed by a ``yes'' entry in this register.

% ============================================================================
% I. Canonical notation
% ============================================================================

\section*{Canonical notation}

\begin{longtable}{
@{}
>{\raggedright\arraybackslash}p{0.18\textwidth}
>{\raggedright\arraybackslash}p{0.18\textwidth}
>{\raggedright\arraybackslash}p{0.24\textwidth}
>{\raggedright\arraybackslash}p{0.16\textwidth}
>{\centering\arraybackslash}p{0.09\textwidth}
>{\raggedleft\arraybackslash}p{0.07\textwidth}
@{}
}

\toprule
\textbf{Notation} & \textbf{Type / signature} & \textbf{Canonical meaning}
& \textbf{Scope} & \textbf{Public} & \textbf{Ref.}\\
\midrule
\endfirsthead
\toprule
\textbf{Notation} & \textbf{Type / signature} & \textbf{Canonical meaning}
& \textbf{Scope} & \textbf{Public} & \textbf{Ref.}\\
\midrule
\endhead
\bottomrule
\endlastfoot

\CanonicalEntry
{\(\Map_{\mathcal C}(X,Y)\)}
{space}
{mapping space in an \(\infty\)-category \(\mathcal C\)}
{global categorical convention}
{yes}
{conv:spectral-infinity-categorical}

\CanonicalEntry
{\(\map_{\mathcal C}(X,Y)\)}
{spectrum}
{mapping spectrum in a stable \(\infty\)-category \(\mathcal C\)}
{global categorical convention}
{yes}
{conv:spectral-infinity-categorical}

\CanonicalEntry
{\(\simeq\)}
{relation / arrow decoration}
{equivalence in an \(\infty\)-category}
{global categorical convention}
{no}
{conv:spectral-infinity-categorical}

\CanonicalEntry
{\(\mathbb U\subset\mathbb V\)}
{Grothendieck universes}
{small-data and ambient-size universes used throughout the chapter}
{size convention}
{no}
{conv:size-spectral-ag}

\CanonicalEntry
{\(\Sh_\tau(\mathcal C)\)}
{\(\infty\)-category}
{\(\tau\)-sheaves of spaces, not hypercompleted unless stated}
{sheaf convention}
{no}
{conv:spectral-sheaves-hyperdescent}

\CanonicalEntry
{\(\Sh_\tau(\mathcal C)^\wedge\)}
{hypercomplete \(\infty\)-topos}
{hypercompletion of the \(\tau\)-sheaf \(\infty\)-topos}
{sheaf convention}
{no}
{conv:spectral-sheaves-hyperdescent}

\CanonicalEntry
{\(\CAlg(\Cat_\infty)\)}
{\(\infty\)-category}
{symmetric monoidal \(\infty\)-categories for the Cartesian monoidal structure on \(\Cat_\infty\)}
{monoidal-category convention}
{no}
{conv:spectral-monoidal-structures-on-categories}

\CanonicalEntry
{\(\CAlg(\Pr^L_{\mathrm{st}})\)}
{\(\infty\)-category}
{presentable stable symmetric monoidal \(\infty\)-categories and colimit-preserving strong symmetric monoidal functors}
{monoidal-category convention}
{yes}
{conv:spectral-monoidal-structures-on-categories}

\CanonicalEntry
{\(\Sp\)}
{presentable stable \(\infty\)-category}
{category of spectra}
{spectral coefficients}
{yes}
{prop:spectra-presentably-stable-monoidal}

\CanonicalEntry
{\(\Sp^\otimes\)}
{object of \(\CAlg(\Pr^L_{\mathrm{st}})\)}
{canonical symmetric monoidal category of spectra}
{spectral coefficients}
{yes}
{prop:spectra-presentably-stable-monoidal}

\CanonicalEntry
{\(\mathbb S\)}
{spectrum}
{sphere spectrum, the tensor unit of \(\Sp^\otimes\)}
{spectral coefficients}
{yes}
{prop:spectra-presentably-stable-monoidal}

\CanonicalEntry
{\(\HHom_{\Sp}(X,-)\)}
{\(\Sp\to\Sp\)}
{right adjoint to \(-\otimes X\)}
{closed spectral monoidal structure}
{yes}
{prop:spectra-closedness}

\CanonicalEntry
{\(\Sp_{\geq n}\)}
{full subcategory of \(\Sp\)}
{spectra with \(\pi_i=0\) for \(i<n\)}
{Postnikov \(t\)-structure}
{yes}
{def:postnikov-t-structure}

\CanonicalEntry
{\(\Sp_{\leq n}\)}
{full subcategory of \(\Sp\)}
{spectra with \(\pi_i=0\) for \(i>n\)}
{Postnikov \(t\)-structure}
{yes}
{def:postnikov-t-structure}

\CanonicalEntry
{\(\tau_{\geq n}^{\Sp}\)}
{Postnikov truncation functor}
{connective truncation at degree \(n\)}
{Postnikov \(t\)-structure on spectra}
{yes}
{def:postnikov-t-structure}

\CanonicalEntry
{\(\tau_{\leq n}^{\Sp}\)}
{Postnikov truncation functor}
{coconnective truncation at degree \(n\)}
{Postnikov \(t\)-structure on spectra}
{yes}
{def:postnikov-t-structure}

\CanonicalEntry
{\(\Sp^\heartsuit\)}
{abelian category}
{heart \(\Sp_{\geq0}\cap\Sp_{\leq0}\)}
{Postnikov \(t\)-structure}
{yes}
{def:discrete-spectrum}

\CanonicalEntry
{\(H_{\mathrm{Ab}}\)}
{\(\operatorname{Ab}\xrightarrow{\simeq}\Sp^\heartsuit\)}
{Eilenberg--Mac Lane equivalence from abelian groups to discrete spectra}
{discrete spectra}
{yes}
{prop:eilenberg-mac-lane-heart}

\CanonicalEntry
{\(H_{\mathrm{CRing}}\)}
{\(\operatorname{CRing}\hookrightarrow\CAlg(\Sp)^{\mathrm{cn}}\)}
{Eilenberg--Mac Lane embedding of ordinary commutative rings as discrete connective commutative ring spectra}
{discrete spectral algebras}
{yes}
{prop:discrete-commutative-ring-spectra}

\CanonicalEntry
{\(\Comm^\otimes\)}
{\(\infty\)-operad}
{commutative \(\infty\)-operad}
{commutative spectral algebra}
{no}
{def:commutative-ring-spectrum}

\CanonicalEntry
{\(\operatorname{Env}(\Comm)^\otimes\)}
{symmetric monoidal envelope}
{symmetric monoidal envelope of the commutative \(\infty\)-operad}
{commutative spectral algebra}
{no}
{rem:spectral-commutative-envelope}

\CanonicalEntry
{\(\CAlg(\Sp)\)}
{\(\infty\)-category}
{commutative ring spectra / \(\EE_\infty\)-rings}
{commutative spectral algebra}
{yes}
{def:commutative-ring-spectrum}

\CanonicalEntry
{\(\operatorname{oblv}_{\CAlg}\)}
{\(\CAlg(\Sp)\to\Sp\)}
{forgetful functor from commutative ring spectra to spectra}
{commutative spectral algebra}
{no}
{prop:calg-spectra-presentable}

\CanonicalEntry
{\(\operatorname{Sym}\)}
{\(\Sp\to\CAlg(\Sp)\)}
{free commutative ring-spectrum functor, left adjoint to \(\operatorname{oblv}_{\CAlg}\)}
{commutative spectral algebra}
{yes}
{def:free-commutative-spectral-algebra}

\CanonicalEntry
{\(\CAlg(\Sp)^{\mathrm{cn}}\)}
{full subcategory of \(\CAlg(\Sp)\)}
{connective commutative ring spectra}
{connective spectral algebra}
{yes}
{def:connective-e-infinity-ring}

\CanonicalEntry
{\(\iota\)}
{\(\CAlg(\Sp)^{\mathrm{cn}}\hookrightarrow\CAlg(\Sp)\)}
{inclusion of connective commutative ring spectra}
{connective spectral algebra}
{no}
{prop:connective-calg-presentable}

\CanonicalEntry
{\(\tau_{\geq0}^{\CAlg}\)}
{\(\CAlg(\Sp)\to\CAlg(\Sp)^{\mathrm{cn}}\)}
{connective-cover right adjoint to \(\iota\)}
{connective spectral algebra}
{yes}
{prop:connective-calg-presentable}

\CanonicalEntry
{\(\pi_0\)}
{\(\CAlg(\Sp)^{\mathrm{cn}}\to\operatorname{CRing}\)}
{degree-zero homotopy / classical truncation on connective commutative ring spectra}
{classical comparison}
{yes}
{prop:pi-zero-h-adjunction}

\CanonicalEntry
{\(\Aff_{\Sp}\)}
{\((\CAlg(\Sp)^{\mathrm{cn}})^{\op}\)}
{\(\infty\)-category of affine spectral schemes}
{affine spectral geometry}
{yes}
{def:affine-spectral-scheme}

\CanonicalEntry
{\(\Spec(A)\)}
{object of \(\Aff_{\Sp}\)}
{affine spectral scheme represented by a connective commutative ring spectrum \(A\)}
{affine spectral geometry}
{yes}
{def:affine-spectral-scheme}

\CanonicalEntry
{\((-)_{\mathrm{cl}}\)}
{affine / prestack / spectral-scheme truncation functor}
{classical truncation, compatible across the three levels constructed in the chapter}
{classical comparison}
{yes}
{thm:global-classical-truncation}

\CanonicalEntry
{\(X_{\mathrm{cl}}\)}
{ordinary affine scheme, classical prestack, or ordinary scheme}
{classical truncation of the spectral object \(X\)}
{classical comparison}
{yes}
{thm:global-classical-truncation}

\CanonicalEntry
{\(|X|\)}
{topological space}
{underlying space of a spectral scheme, defined as \(|X_{\mathrm{cl}}|\)}
{spectral schemes}
{yes}
{def:underlying-space-spectral-scheme}

\CanonicalEntry
{\(\Mod_A\)}
{presentable stable symmetric monoidal \(\infty\)-category}
{category of \(A\)-module spectra}
{modules}
{yes}
{def:module-spectrum-over-a}

\CanonicalEntry
{\(-\otimes_A-\)}
{\(\Mod_A\times\Mod_A\to\Mod_A\)}
{relative tensor product of \(A\)-modules}
{modules}
{yes}
{con:spectral-relative-tensor-product}

\CanonicalEntry
{\(\HHom_A(M,-)\)}
{\(\Mod_A\to\Mod_A\)}
{right adjoint to \(-\otimes_A M\)}
{modules}
{yes}
{prop:module-spectra-closedness}

\CanonicalEntry
{\(f^*\)}
{\(\Mod_A\to\Mod_B\), for \(f:A\to B\)}
{extension of scalars \(-\otimes_A B\); globally, quasi-coherent pullback}
{base change / reindexing}
{yes}
{def:spectral-extension-restriction-scalars}

\CanonicalEntry
{\(f_*\)}
{\(\Mod_B\to\Mod_A\), for \(f:A\to B\)}
{restriction of scalars; globally, right adjoint quasi-coherent direct image}
{base change / reindexing}
{yes}
{def:spectral-extension-restriction-scalars}

\CanonicalEntry
{\(\mathbf{Mod}_{\Sp}\)}
{\(\CAlg(\Sp)\to\CAlg(\Pr^L_{\mathrm{st}})\)}
{operadic module-category functor, with transition functors extension of scalars}
{indexed coefficients}
{yes}
{prop:spectral-module-functoriality}

\CanonicalEntry
{\(\CAlg(\Mod_A)\)}
{\(\infty\)-category}
{commutative algebra objects in \(A\)-modules, canonically \(\CAlg(\Sp)_{A/}\)}
{relative spectral algebra}
{yes}
{prop:relative-commutative-spectral-algebras}

\CanonicalEntry
{\((\Mod_A)_{\geq0}\)}
{full subcategory of \(\Mod_A\)}
{\(A\)-modules with connective underlying spectrum}
{module Postnikov structure}
{yes}
{def:module-postnikov-t-structure}

\CanonicalEntry
{\((\Mod_A)_{\leq0}\)}
{full subcategory of \(\Mod_A\)}
{\(A\)-modules with coconnective underlying spectrum}
{module Postnikov structure}
{yes}
{def:module-postnikov-t-structure}

\CanonicalEntry
{\(\Mod_A^\heartsuit\)}
{abelian category}
{heart of the module Postnikov \(t\)-structure, naturally \(\operatorname{Mod}_{\pi_0(A)}\)}
{module Postnikov structure}
{yes}
{prop:heart-module-category}

\CanonicalEntry
{\(\operatorname{Sym}_A\)}
{\(\Mod_A\to\CAlg(\Mod_A)\)}
{free commutative \(A\)-algebra functor}
{relative spectral algebra}
{yes}
{def:finite-cell-spectral-algebra}

\CanonicalEntry
{\(\operatorname{Res}_A^B\)}
{\(\Mod_B\to\Mod_A\)}
{restriction of scalars, used in the affine direct-image formula}
{affine direct image}
{no}
{prop:affine-direct-image-restriction-scalars}

\CanonicalEntry
{\(\PSh(\Aff_{\Sp})\)}
{\(\Fun((\Aff_{\Sp})^{\op},\mathcal S)\)}
{\(\infty\)-category of spectral prestacks}
{prestack geometry}
{yes}
{def:spectral-prestacks-affine-stage}

\CanonicalEntry
{\(h_A\)}
{\((\Aff_{\Sp})^{\op}\to\mathcal S\)}
{functor of points of \(\Spec(A)\)}
{prestack geometry}
{yes}
{def:spectral-functor-of-points}

\CanonicalEntry
{\(h\)}
{\(\Aff_{\Sp}\hookrightarrow\PSh(\Aff_{\Sp})\)}
{spectral Yoneda embedding}
{prestack geometry}
{yes}
{prop:spectral-yoneda-embedding}

\CanonicalEntry
{\(\delta\)}
{\(\operatorname{AffSch}\hookrightarrow\Aff_{\Sp}\)}
{fully faithful discrete affine enhancement \(\Spec(R)\mapsto\Spec(HR)\)}
{classical comparison of prestacks}
{yes}
{def:classical-truncation-prestack}

\CanonicalEntry
{\(\delta_!\)}
{\(\PSh(\operatorname{AffSch})\to\PSh(\Aff_{\Sp})\)}
{left Kan extension along \(\delta^{\op}\)}
{classical comparison of prestacks}
{yes}
{def:classical-truncation-prestack}

\CanonicalEntry
{\(\delta^*\)}
{\(\PSh(\Aff_{\Sp})\to\PSh(\operatorname{AffSch})\)}
{restriction to discrete affine spectral schemes; prestack classical truncation}
{classical comparison of prestacks}
{yes}
{def:classical-truncation-prestack}

\CanonicalEntry
{\(\QCoh_{\mathrm{aff}}\)}
{\((\Aff_{\Sp})^{\op}\to\CAlg(\Pr^L_{\mathrm{st}})\)}
{affine quasi-coherent coefficient system, \(\Spec(A)\mapsto\Mod_A\)}
{quasi-coherent coefficients}
{yes}
{prop:affine-qcoh-coefficient-system}

\CanonicalEntry
{\(\operatorname{AffPt}(X)\)}
{\(\infty\)-category}
{category of affine points \(\Spec(A)\to X\) of a spectral prestack}
{quasi-coherent coefficients}
{yes}
{def:affine-points-prestack}

\CanonicalEntry
{\(\QCoh(X)\)}
{object of \(\CAlg(\Pr^L_{\mathrm{st}})\)}
{quasi-coherent coefficient category, the affine right Kan extension of \(\QCoh_{\mathrm{aff}}\)}
{quasi-coherent coefficients}
{yes}
{def:qcoh-spectral-prestack}

\CanonicalEntry
{\(\mathcal O_X\)}
{object of \(\QCoh(X)\)}
{tensor unit, called the structure sheaf}
{quasi-coherent coefficients}
{yes}
{def:structure-sheaf-prestack}

\CanonicalEntry
{\(\Gamma(X,\mathcal F)\)}
{spectrum}
{global sections \(\map_{\QCoh(X)}(\mathcal O_X,\mathcal F)\)}
{quasi-coherent coefficients}
{yes}
{def:global-sections-qcoh}

\CanonicalEntry
{\(A[f^{-1}]\)}
{commutative \(A\)-algebra}
{spectral localization initial among \(A\)-algebras in which \(f\in\pi_0(A)\) is invertible}
{principal opens}
{yes}
{con:spectral-localization}

\CanonicalEntry
{\(D_A^{\mathrm{sp}}(f)\)}
{affine spectral scheme over \(\Spec(A)\)}
{principal spectral open \(\Spec(A[f^{-1}])\)}
{principal opens}
{yes}
{prop:principal-spectral-open}

\CanonicalEntry
{\(\tau_{\mathrm{Zar}}\)}
{Grothendieck topology on \(\Aff_{\Sp}\)}
{spectral Zariski topology generated by affine spectral Zariski covers}
{Zariski geometry}
{yes}
{def:spectral-zariski-topology}

\CanonicalEntry
{\(\operatorname{Stk}^{\mathrm{sp}}_{\mathrm{Zar}}\)}
{\(\Sh_{\tau_{\mathrm{Zar}}}(\Aff_{\Sp})\)}
{\(\infty\)-category of spectral Zariski stacks}
{global spectral geometry}
{yes}
{def:spectral-zariski-stack}

\CanonicalEntry
{\(a^{\mathrm{sp}}_{\mathrm{Zar}}\)}
{\(\PSh(\Aff_{\Sp})\to\operatorname{Stk}^{\mathrm{sp}}_{\mathrm{Zar}}\)}
{spectral Zariski sheafification}
{global spectral geometry}
{yes}
{con:spectral-zariski-sheafification}

\CanonicalEntry
{\(\operatorname{Sch}^{\mathrm{sp}}\)}
{full subcategory of \(\operatorname{Stk}^{\mathrm{sp}}_{\mathrm{Zar}}\)}
{spectral schemes, i.e. spectral Zariski stacks with affine open atlases}
{global spectral geometry}
{yes}
{def:spectral-scheme}

\CanonicalEntry
{\(\operatorname{Sch}^{\mathrm{sp}}_{\mathrm{str}}\)}
{\(\infty\)-category}
{connective \(0\)-localic structured spectral schemes used in the comparison theorem}
{model comparison}
{no}
{thm:spectral-scheme-model-comparison}

\CanonicalEntry
{\(\Phi\)}
{\(\operatorname{Sch}^{\mathrm{sp}}_{\mathrm{str}}\to\operatorname{Stk}^{\mathrm{sp}}_{\mathrm{Zar}}\)}
{functor-of-points comparison functor}
{model comparison}
{no}
{thm:spectral-scheme-model-comparison}

\CanonicalEntry
{\(d\)}
{\(\Sh_{\mathrm{Zar}}(\operatorname{AffSch})\to\operatorname{Stk}^{\mathrm{sp}}_{\mathrm{Zar}}\)}
{discrete extension \(a^{\mathrm{sp}}_{\mathrm{Zar}}\circ\delta_!\), left adjoint to \(\delta^*\)}
{classical comparison}
{yes}
{con:discrete-extension-zariski-sheaves}

\CanonicalEntry
{\((-)^{\mathrm{disc}}\)}
{\(\operatorname{Sch}\hookrightarrow\operatorname{Sch}^{\mathrm{sp}}\)}
{fully faithful discrete spectral enhancement of ordinary schemes}
{classical comparison}
{yes}
{thm:ordinary-schemes-discrete-enhancement}

\CanonicalEntry
{\(i_X\)}
{\(X_{\mathrm{cl}}^{\mathrm{disc}}\to X\)}
{natural classical comparison closed immersion from the discrete enhancement of the classical truncation}
{classical comparison}
{yes}
{cor:classical-comparison-morphism}

\CanonicalEntry
{\(\operatorname{AffOp}(X)\)}
{\(\infty\)-category}
{affine open spectral subschemes of \(X\)}
{global quasi-coherent coefficients}
{yes}
{lem:qcoh-affine-open-basis}

\CanonicalEntry
{\(\QCoh(X)_{\geq0}\)}
{full subcategory of \(\QCoh(X)\)}
{quasi-coherent modules connective on every affine open}
{global Postnikov structure}
{yes}
{thm:qcoh-spectral-scheme-t-structure}

\CanonicalEntry
{\(\QCoh(X)_{\leq0}\)}
{full subcategory of \(\QCoh(X)\)}
{quasi-coherent modules coconnective on every affine open}
{global Postnikov structure}
{yes}
{thm:qcoh-spectral-scheme-t-structure}

\CanonicalEntry
{\(\QCoh(X)^\heartsuit\)}
{abelian category}
{heart of the global Postnikov \(t\)-structure}
{global Postnikov structure}
{yes}
{thm:qcoh-heart-classical-truncation}

\CanonicalEntry
{\(\operatorname{QCAlg}(X)\)}
{\(\CAlg(\QCoh(X))\)}
{quasi-coherent commutative algebra objects on \(X\)}
{relative affine geometry}
{yes}
{def:quasi-coherent-commutative-algebra}

\CanonicalEntry
{\(\operatorname{oblv}^{\mathrm{alg}}_X\)}
{\(\operatorname{QCAlg}(X)\to\QCoh(X)\)}
{forgetful functor from quasi-coherent commutative algebras}
{relative affine geometry}
{no}
{def:quasi-coherent-commutative-algebra}

\CanonicalEntry
{\(\operatorname{QCAlg}(X)^{\mathrm{cn}}\)}
{full subcategory of \(\operatorname{QCAlg}(X)\)}
{algebra objects with connective underlying quasi-coherent module}
{relative affine geometry}
{yes}
{def:quasi-coherent-commutative-algebra}

\CanonicalEntry
{\(\Spec_X\)}
{\(\operatorname{QCAlg}(X)^{\mathrm{cn},\op}\to\operatorname{Sch}^{\mathrm{sp}}_{/X}\)}
{relative spectrum functor classifying affine morphisms over \(X\)}
{relative affine geometry}
{yes}
{thm:relative-spectrum-construction}

\CanonicalEntry
{\(X\setminus Z\)}
{open spectral subscheme of \(X\)}
{open complement of a closed immersion \(Z\to X\)}
{closed geometry}
{yes}
{thm:complement-closed-immersion}

\CanonicalEntry
{\(\operatorname{Sym}_{\mathcal O_X}\)}
{\(\QCoh(X)\to\operatorname{QCAlg}(X)\)}
{relative free commutative algebra functor}
{linear geometry}
{yes}
{prop:relative-free-commutative-algebra}

\CanonicalEntry
{\(\mathbb L_X(\mathcal E)\)}
{spectral scheme over \(X\)}
{relative linear space \(\Spec_X\operatorname{Sym}_{\mathcal O_X}(\mathcal E)\)}
{linear geometry}
{yes}
{def:linear-spectral-space}

\CanonicalEntry
{\(\mathbf V_X(\mathcal E)\)}
{spectral scheme over \(X\)}
{section total space \(\mathbb L_X(\mathcal E^\vee)\) for finite locally free \(\mathcal E\)}
{linear geometry}
{yes}
{def:vector-bundle-dual-convention}

\CanonicalEntry
{\(\mathbb A_X^n\)}
{spectral scheme over \(X\)}
{brave new affine \(n\)-space \(\mathbb L_X(\mathcal O_X^{\oplus n})\)}
{distinguished affine objects}
{yes}
{def:brave-new-affine-space}

\CanonicalEntry
{\(\mathbb A_X^1\)}
{spectral scheme over \(X\)}
{brave new affine line}
{distinguished affine objects}
{yes}
{def:brave-new-affine-space}

\CanonicalEntry
{\(0,1\)}
{\(X\rightrightarrows\mathbb A_X^1\)}
{zero and unit sections induced by the zero and identity endomorphisms of \(\mathcal O_X\)}
{distinguished affine objects}
{yes}
{con:zero-unit-sections-affine-line}

\CanonicalEntry
{\(\mathbb G_{m,X}\)}
{commutative group object over \(X\)}
{open complement \(\mathbb A_X^1\setminus0(X)\)}
{distinguished affine objects}
{yes}
{def:spectral-multiplicative-group}

\CanonicalEntry
{\(A\{t\}\)}
{commutative \(A\)-algebra}
{brave new polynomial algebra \(\operatorname{Sym}_A(A)\)}
{affine line coordinates}
{yes}
{prop:spectral-multiplicative-group-structure}

\CanonicalEntry
{\(R\{t\}\)}
{commutative \(HR\)-algebra}
{\(\operatorname{Sym}_{HR}(HR)\), used for the homotopy calculation over an ordinary ring \(R\)}
{affine line coordinates}
{yes}
{prop:affine-line-genuinely-spectral}

\end{longtable}

% ============================================================================
% II. Canonical terminology
% ============================================================================

\section*{Canonical terminology}

This section controls technical terms whose meanings must remain stable across later
chapters.  It is intentionally definitional rather than encyclopedic.

\begin{longtable}{
@{}
>{\raggedright\arraybackslash}p{0.24\textwidth}
>{\raggedright\arraybackslash}p{0.50\textwidth}
>{\raggedright\arraybackslash}p{0.14\textwidth}
>{\raggedleft\arraybackslash}p{0.08\textwidth}
@{}
}

\toprule
\textbf{Term} & \textbf{Canonical use} & \textbf{Scope} & \textbf{Ref.}\\
\midrule
\endfirsthead
\toprule
\textbf{Term} & \textbf{Canonical use} & \textbf{Scope} & \textbf{Ref.}\\
\midrule
\endhead
\bottomrule
\endlastfoot

\CanonicalTerm
{spectral prestack}
{a space-valued presheaf on \(\Aff_{\Sp}\)}
{prestack geometry}
{def:spectral-prestacks-affine-stage}

\CanonicalTerm
{affine spectral scheme}
{an object \(\Spec(A)\) of \(\Aff_{\Sp}=(\CAlg(\Sp)^{\mathrm{cn}})^{\op}\)}
{affine spectral geometry}
{def:affine-spectral-scheme}

\CanonicalTerm
{commutative ring spectrum / \(\EE_\infty\)-ring}
{a commutative algebra object in spectra}
{commutative spectral algebra}
{def:commutative-ring-spectrum}

\CanonicalTerm
{connective \(\EE_\infty\)-ring}
{a commutative ring spectrum whose underlying spectrum is connective}
{connective spectral algebra}
{def:connective-e-infinity-ring}

\CanonicalTerm
{discrete spectrum}
{an object of \(\Sp_{\geq0}\cap\Sp_{\leq0}\)}
{Postnikov theory}
{def:discrete-spectrum}

\CanonicalTerm
{module spectrum}
{an object of \(\Mod_A\) for a connective commutative ring spectrum \(A\)}
{modules}
{def:module-spectrum-over-a}

\CanonicalTerm
{flat module spectrum}
{a connective \(A\)-module \(M\) for which \(-\otimes_A M\) is \(t\)-exact}
{geometric exactness}
{def:flat-module-spectrum}

\CanonicalTerm
{flat spectral algebra morphism}
{a map \(A\to B\) for which \(B\) is flat as an \(A\)-module}
{geometric exactness}
{def:flat-spectral-algebra-morphism}

\CanonicalTerm
{locally of finite presentation}
{for \(A\to B\), compactness of \(B\) in \(\CAlg(\Sp)_{A/}\)}
{finiteness}
{def:finitely-presented-spectral-algebra}

\CanonicalTerm
{finite cell algebra}
{a connective commutative \(A\)-algebra in the finite-colimit closure of the finitely generated free commutative \(A\)-algebras}
{finiteness}
{def:finite-cell-spectral-algebra}

\CanonicalTerm
{perfect module}
{an \(A\)-module in the thick subcategory generated by the tensor unit \(A\)}
{module finiteness}
{def:perfect-module-spectrum}

\CanonicalTerm
{Tor-amplitude in \([a,b]\)}
{the condition that tensoring with every discrete \(A\)-module has homotopy only in degrees \(a,\ldots,b\)}
{module finiteness}
{def:tor-amplitude}

\CanonicalTerm
{finite projective module}
{a retract of a finite free \(A\)-module}
{module finiteness}
{def:finite-projective-module-spectrum}

\CanonicalTerm
{objectwise connective quasi-coherent module}
{a quasi-coherent module whose pullback to every affine point is connective}
{prestack coefficients}
{def:objectwise-connective-qcoh}

\CanonicalTerm
{affine open immersion}
{an affine map represented by a flat, locally finitely presented algebra map with idempotent multiplication \(B\otimes_A B\simeq B\)}
{affine Zariski geometry}
{def:affine-open-immersion-spectral}

\CanonicalTerm
{affine spectral Zariski covering family}
{a family of affine open immersions whose classical opens cover \(\Spec\pi_0(A)\)}
{affine Zariski geometry}
{def:affine-zariski-cover}

\CanonicalTerm
{spectral Zariski topology}
{the Grothendieck topology on \(\Aff_{\Sp}\) generated by affine spectral Zariski covers}
{Zariski geometry}
{def:spectral-zariski-topology}

\CanonicalTerm
{descendable morphism}
{a map \(A\to B\) for which \(B\) generates \(\Mod_A\) as a thick tensor ideal}
{descent}
{def:spectral-descendable-morphism}

\CanonicalTerm
{spectral Zariski stack}
{a \(\tau_{\mathrm{Zar}}\)-sheaf of spaces on \(\Aff_{\Sp}\)}
{global spectral geometry}
{def:spectral-zariski-stack}

\CanonicalTerm
{open immersion}
{a representable morphism of spectral Zariski stacks whose affine base changes are covered by affine spectral open immersions}
{global spectral geometry}
{def:open-immersion-spectral-stacks}

\CanonicalTerm
{spectral scheme}
{a spectral Zariski stack admitting an effective epimorphic affine open atlas}
{global spectral geometry}
{def:spectral-scheme}

\CanonicalTerm
{classical truncation}
{restriction to discrete affine spectral schemes, shown on spectral schemes to land in ordinary schemes}
{classical comparison}
{thm:global-classical-truncation}

\CanonicalTerm
{discrete enhancement}
{the fully faithful spectral enhancement of an ordinary scheme obtained from Eilenberg--Mac Lane affines and Zariski gluing}
{classical comparison}
{thm:ordinary-schemes-discrete-enhancement}

\CanonicalTerm
{quasi-compact spectral scheme}
{a spectral scheme admitting a finite affine open cover}
{global finiteness}
{def:quasi-compact-spectral-scheme}

\CanonicalTerm
{quasi-separated morphism}
{a morphism whose diagonal is quasi-compact}
{global finiteness}
{def:quasi-separated-spectral-scheme}

\CanonicalTerm
{qcqs spectral scheme}
{a spectral scheme that is both quasi-compact and quasi-separated}
{global finiteness}
{def:quasi-separated-spectral-scheme}

\CanonicalTerm
{flat morphism of spectral schemes}
{a morphism whose affine-local coordinate algebra maps are flat}
{global geometric exactness}
{def:flat-spectral-scheme-morphism}

\CanonicalTerm
{flat quasi-coherent module}
{a quasi-coherent module whose restriction to every affine open is flat}
{global geometric exactness}
{def:flat-spectral-scheme-morphism}

\CanonicalTerm
{affine morphism of spectral schemes}
{a morphism whose pullback over every affine open of the target is affine}
{relative affine geometry}
{def:affine-morphism-spectral-schemes}

\CanonicalTerm
{quasi-coherent commutative algebra}
{a commutative algebra object of \(\QCoh(X)\)}
{relative affine geometry}
{def:quasi-coherent-commutative-algebra}

\CanonicalTerm
{relative spectrum}
{the spectral scheme \(\Spec_X(\mathcal A)\) representing algebra maps \(f^*\mathcal A\to\mathcal O_Y\)}
{relative affine geometry}
{thm:relative-spectrum-construction}

\CanonicalTerm
{affine closed immersion}
{an affine spectral map \(\Spec(B)\to\Spec(A)\) for which \(\pi_0(A)\to\pi_0(B)\) is surjective}
{closed geometry}
{def:affine-closed-immersion-spectral}

\CanonicalTerm
{closed immersion}
{a morphism of spectral schemes whose pullback to every affine open is an affine closed immersion}
{closed geometry}
{def:closed-immersion-spectral-schemes}

\CanonicalTerm
{nil-immersion}
{a closed immersion inducing an equivalence on reduced classical schemes}
{closed geometry}
{def:nil-immersion-spectral}

\CanonicalTerm
{open complement}
{the open \(X\setminus Z\) characterized by vanishing of the base-changed closed subscheme}
{closed geometry}
{thm:complement-closed-immersion}

\CanonicalTerm
{locally free module of rank \(n\)}
{an \(A\)-module that becomes \(A_i^{\oplus n}\) on a spectral Zariski cover}
{linear geometry}
{def:locally-free-affine-module}

\CanonicalTerm
{vector bundle of rank \(n\)}
{a quasi-coherent module locally isomorphic to \(\mathcal O^{\oplus n}\)}
{linear geometry}
{def:locally-free-qcoh}

\CanonicalTerm
{relative linear space}
{\(\mathbb L_X(\mathcal E)=\Spec_X\operatorname{Sym}_{\mathcal O_X}(\mathcal E)\) for connective \(\mathcal E\)}
{linear geometry}
{def:linear-spectral-space}

\CanonicalTerm
{section total space}
{\(\mathbf V_X(\mathcal E)=\mathbb L_X(\mathcal E^\vee)\) for finite locally free \(\mathcal E\), representing sections of \(\mathcal E\)}
{linear geometry}
{def:vector-bundle-dual-convention}

\CanonicalTerm
{brave new affine space}
{\(\mathbb A_X^n=\mathbb L_X(\mathcal O_X^{\oplus n})\)}
{distinguished affine objects}
{def:brave-new-affine-space}

\CanonicalTerm
{brave new affine line}
{the case \(\mathbb A_X^1\) of brave new affine space}
{distinguished affine objects}
{def:brave-new-affine-space}

\CanonicalTerm
{spectral multiplicative group}
{the open complement \(\mathbb G_{m,X}=\mathbb A_X^1\setminus0(X)\)}
{distinguished affine objects}
{def:spectral-multiplicative-group}

\end{longtable}

% ============================================================================
% III. Canonical identities and status checks
% ============================================================================

\section*{Canonical identities and status checks}

The entries in this section are not additional meanings for the symbols above.  They are
equations, comparison maps, and status facts that later chapters may use to audit whether
the registered notation is being transported correctly.

\begin{longtable}{
@{}
>{\raggedright\arraybackslash}p{0.22\textwidth}
>{\raggedright\arraybackslash}p{0.61\textwidth}
>{\raggedleft\arraybackslash}p{0.13\textwidth}
@{}
}

\toprule
\textbf{Check} & \textbf{Canonical identity or status} & \textbf{Ref.}\\
\midrule
\endfirsthead
\toprule
\textbf{Check} & \textbf{Canonical identity or status} & \textbf{Ref.}\\
\midrule
\endhead
\bottomrule
\endlastfoot

\CanonicalLaw
{Infinity-categorical convention}
{All categories are \(\infty\)-categories unless explicitly stated; \(\Map\) denotes mapping spaces and \(\map\) mapping spectra in stable categories.}
{conv:spectral-infinity-categorical}

\CanonicalLaw
{Higher coherence}
{Functorial, adjoint, coCartesian, operadic, Kan-extension, limit, and descent constructions carry the coherent higher structure supplied by those constructions; no extra coherence data are imposed.}
{conv:spectral-higher-coherence}

\CanonicalLaw
{Sheaf convention}
{\(\Sh_\tau(\mathcal C)\) means sheaves of spaces without hypercompletion; hypercomplete sheaves are written \(\Sh_\tau(\mathcal C)^\wedge\).}
{conv:spectral-sheaves-hyperdescent}

\CanonicalLaw
{Affine variance}
{\(\Aff_{\Sp}=(\CAlg(\Sp)^{\mathrm{cn}})^{\op}\); hence \(\Map_{\Aff_{\Sp}}(\Spec B,\Spec A)\simeq\Map_{\CAlg(\Sp)}(A,B)\).}
{prop:affine-spectral-mapping-spaces}

\CanonicalLaw
{Affine fibre products}
{\(\Spec(B)\times_{\Spec(A)}\Spec(C)\simeq\Spec(B\otimes_A C)\).}
{prop:affine-spectral-fibre-products}

\CanonicalLaw
{Classical truncation adjunction}
{\(\pi_0\dashv H_{\mathrm{CRing}}\) on connective commutative ring spectra; the unit is \(A\to H\pi_0(A)\).}
{prop:pi-zero-h-adjunction}

\CanonicalLaw
{Affine classical truncation}
{\((\Spec A)_{\mathrm{cl}}\simeq\Spec\pi_0(A)\), and affine classical truncation preserves fibre products.}
{rem:classical-truncation-affine-left-exact}

\CanonicalLaw
{Module base change}
{For \(f:A\to B\), \(f^*=-\otimes_A B\dashv f_*\); \(f^*\) is colimit-preserving and strong symmetric monoidal.}
{prop:spectral-base-change-monoidal}

\CanonicalLaw
{Stable exactness versus flatness}
{Extension of scalars is always exact as a stable functor; geometric flatness is the stronger requirement of \(t\)-exactness.}
{rem:stable-exactness-not-flatness}

\CanonicalLaw
{Spectral Lazard criterion}
{For connective \(M\), flatness is equivalent to ordinary flatness of \(\pi_0(M)\) plus \(\pi_n(M)\cong\pi_n(A)\otimes_{\pi_0(A)}\pi_0(M)\), and also to being a filtered colimit of finite free modules.}
{thm:spectral-lazard}

\CanonicalLaw
{Finite-presentation criterion}
{A connective commutative \(A\)-algebra is locally of finite presentation iff it is a retract of a finite cell \(A\)-algebra.}
{prop:finite-presentation-compactness-criterion}

\CanonicalLaw
{Finite projective criterion}
{Finite projective, perfect of Tor-amplitude \([0,0]\), and perfect-connective-flat are equivalent conditions for an \(A\)-module.}
{prop:finite-projective-characterization}

\CanonicalLaw
{Prestack classical truncation}
{For a spectral prestack \(X\), \(X_{\mathrm{cl}}=\delta^*X\); restriction preserves all limits and sends \(\Spec A\) to \(\Spec\pi_0(A)\).}
{prop:prestack-classical-truncation}

\CanonicalLaw
{QCoh right Kan extension}
{\(\QCoh\) is the right Kan extension of \(\QCoh_{\mathrm{aff}}\) along opposite Yoneda, with \(\QCoh(X)=\varprojlim_{(\Spec A\to X)\in\operatorname{AffPt}(X)^{\op}}\Mod_A\).}
{thm:qcoh-affine-right-kan-extension}

\CanonicalLaw
{Affine QCoh recovery}
{\(\QCoh(\Spec A)\simeq\Mod_A\) symmetrically monoidally.}
{cor:qcoh-recovery-affines}

\CanonicalLaw
{QCoh variance on colimits}
{\(\QCoh(\varinjlim_\alpha X_\alpha)\simeq\varprojlim_\alpha\QCoh(X_\alpha)\).}
{cor:qcoh-colimits-to-limits}

\CanonicalLaw
{Principal localization}
{\(\pi_n(A[f^{-1}])\cong\pi_n(A)[f^{-1}]\); in particular \(A\to A[f^{-1}]\) is flat.}
{prop:homotopy-groups-localization}

\CanonicalLaw
{Affine-open recognition}
{\(\Spec(B)\to\Spec(A)\) is an affine spectral open immersion iff \(A\to B\) is flat and \(\Spec\pi_0(B)\to\Spec\pi_0(A)\) is an ordinary open immersion.}
{thm:classical-characterization-affine-opens}

\CanonicalLaw
{Finite principal refinements}
{Every affine spectral Zariski cover refines to finitely many principal opens \(D_A^{\mathrm{sp}}(f_j)\) with the \(f_j\) generating the unit ideal.}
{prop:finite-principal-refinement}

\CanonicalLaw
{Zariski joint conservativity}
{An affine open family covers iff the pullback functors \(-\otimes_A B_i\) are jointly conservative on \(\Mod_A\).}
{prop:zariski-cover-joint-conservativity}

\CanonicalLaw
{Finite-cover descendability}
{For a finite affine spectral Zariski cover, \(A\to\prod_iB_i\) is descendable.}
{prop:finite-zariski-covers-descendable}

\CanonicalLaw
{Amitsur convergence}
{For descendable \(A\to B\), \(A\simeq\operatorname*{Tot}_{[n]\in\Delta}B^{\otimes_A(n+1)}\).}
{thm:spectral-amitsur-convergence}

\CanonicalLaw
{Affine module descent}
{For descendable \(A\to B\), \(\Mod_A\simeq\operatorname*{Tot}_{[n]\in\Delta}\Mod_{B^{\otimes_A(n+1)}}\).}
{thm:spectral-affine-module-descent}

\CanonicalLaw
{Affine algebra descent}
{For descendable \(A\to B\), commutative \(A\)-algebras satisfy the corresponding Amitsur totalization equivalence.}
{thm:spectral-affine-algebra-descent}

\CanonicalLaw
{Zariski subcanonicity}
{Every representable spectral prestack is a \(\tau_{\mathrm{Zar}}\)-sheaf.}
{thm:spectral-zariski-subcanonical}

\CanonicalLaw
{Global finite limits}
{Spectral schemes are closed under fibre products and therefore admit finite limits computed in spectral Zariski stacks.}
{cor:spectral-schemes-finite-limits}

\CanonicalLaw
{Model comparison}
{The functor of points gives an equivalence from connective \(0\)-localic structured spectral schemes to the spectral schemes defined in this chapter.}
{thm:spectral-scheme-model-comparison}

\CanonicalLaw
{Discrete extension adjunction}
{\(d=a^{\mathrm{sp}}_{\mathrm{Zar}}\circ\delta_!\), \(d\dashv\delta^*\), and \(d(\Spec R)\simeq\Spec(HR)\).}
{prop:discrete-extension-affine-calculation}

\CanonicalLaw
{Discrete retraction}
{For every ordinary scheme \(X\), \((X^{\mathrm{disc}})_{\mathrm{cl}}\simeq X\).}
{thm:global-classical-truncation}

\CanonicalLaw
{Classical open control}
{Open spectral subschemes of \(X\) and open subschemes of \(X_{\mathrm{cl}}\) form equivalent posets.}
{prop:opens-classical-truncation}

\CanonicalLaw
{Classical affineness detection}
{A spectral scheme is affine iff its classical truncation is affine.}
{thm:affineness-detected-classically}

\CanonicalLaw
{Classical qcqs detection}
{Quasi-compactness, quasi-separatedness, and qcqs are detected by classical truncation.}
{prop:classical-detection-qcqs}

\CanonicalLaw
{Atlas computation of QCoh}
{For an affine Zariski atlas \(U\to X\), \(\QCoh(X)\simeq\operatorname*{Tot}_{[n]\in\Delta}\QCoh(U_n)\).}
{thm:qcoh-atlas-computation}

\CanonicalLaw
{Affine-open basis for QCoh}
{\(\QCoh(X)\simeq\varprojlim_{V\in\operatorname{AffOp}(X)^{\op}}\QCoh(V)\).}
{lem:qcoh-affine-open-basis}

\CanonicalLaw
{Global Postnikov structure}
{Connectivity and coconnectivity on all affine opens define a complete accessible \(t\)-structure on \(\QCoh(X)\); arbitrary pullback is right \(t\)-exact and flat pullback is \(t\)-exact.}
{thm:qcoh-spectral-scheme-t-structure}

\CanonicalLaw
{Classical heart}
{\(\QCoh(X)^\heartsuit\simeq\operatorname{QCoh}_{\mathrm{ab}}(X_{\mathrm{cl}})\).}
{thm:qcoh-heart-classical-truncation}

\CanonicalLaw
{QCAlg descent}
{Quasi-coherent commutative algebras and their connective subcategories satisfy Zariski descent.}
{prop:qcalg-zariski-descent}

\CanonicalLaw
{Relative-spectrum universal property}
{\(\Map_X(Y,\Spec_X\mathcal A)\simeq\Map_{\operatorname{QCAlg}(Y)}(f^*\mathcal A,\mathcal O_Y)\).}
{thm:relative-spectrum-construction}

\CanonicalLaw
{Relative-spectrum base change}
{\(\Spec_X(\mathcal A)\times_XY\simeq\Spec_Y(f^*\mathcal A)\).}
{prop:relative-spectrum-base-change}

\CanonicalLaw
{Affine-morphism classification}
{\(\operatorname{QCAlg}(X)^{\mathrm{cn},\op}\simeq(\operatorname{Sch}^{\mathrm{sp}}_{/X})_{\mathrm{aff}}\), with inverse \(p\mapsto p_*\mathcal O_Y\).}
{thm:classification-affine-morphisms}

\CanonicalLaw
{Closed-immersion detection}
{A morphism of spectral schemes is a closed immersion iff its classical truncation is a closed immersion.}
{prop:closed-immersions-classical-detection}

\CanonicalLaw
{Classical comparison is closed}
{\(i_X:X_{\mathrm{cl}}^{\mathrm{disc}}\to X\) is a nil-immersion.}
{cor:classical-truncation-nil-immersion}

\CanonicalLaw
{Complement universal property}
{A map \(T\to X\) factors through \(X\setminus Z\) iff \(T\times_XZ\simeq\varnothing\); classical truncation gives the ordinary open complement.}
{thm:complement-closed-immersion}

\CanonicalLaw
{Closed pushouts}
{Pushouts along two closed immersions exist in spectral schemes, the canonical maps are closed, and the construction is stable under arbitrary base change.}
{thm:closed-pushouts-spectral-schemes}

\CanonicalLaw
{Vector bundle criterion}
{Finite projective, Zariski-locally finite free of finite rank, and perfect-connective-flat are equivalent for \(A\)-modules.}
{thm:finite-projective-locally-free}

\CanonicalLaw
{Relative free algebra base change}
{\(f^*\operatorname{Sym}_{\mathcal O_X}(\mathcal E)\simeq\operatorname{Sym}_{\mathcal O_Y}(f^*\mathcal E)\).}
{prop:relative-free-commutative-algebra}

\CanonicalLaw
{Linear-space functor of points}
{\(\Map_X(T,\mathbb L_X(\mathcal E))\simeq\Map_{\QCoh(T)}(f^*\mathcal E,\mathcal O_T)\).}
{prop:linear-space-functor-points}

\CanonicalLaw
{Section-total-space functor of points}
{For finite locally free \(\mathcal E\), \(\Map_X(T,\mathbf V_X(\mathcal E))\simeq\Map_{\QCoh(T)}(\mathcal O_T,f^*\mathcal E)\).}
{prop:vector-bundle-total-space-sections}

\CanonicalLaw
{Affine-space base change}
{\(\mathbb A_X^n\times_XY\simeq\mathbb A_Y^n\), and \((\mathbb A_X^n)_{\mathrm{cl}}\simeq\mathbf A_{X_{\mathrm{cl}}}^n\).}
{prop:classical-truncation-affine-space}

\CanonicalLaw
{Spectral multiplicative group}
{\(\Map_X(T,\mathbb G_{m,X})\simeq\operatorname{Aut}_{\QCoh(T)}(\mathcal O_T)\); the construction commutes with base change.}
{prop:spectral-multiplicative-group-structure}

\CanonicalLaw
{Brave polynomial homotopy}
{For ordinary \(R\), \(R\{t\}\simeq\bigoplus_{n\geq0}HR\otimes\Sigma^\infty_+B\Sigma_n\), with \(\pi_0(R\{t\})\cong R[t]\).}
{prop:affine-line-genuinely-spectral}

\end{longtable}

% ============================================================================
% IV. Intentional notation overloads
% ============================================================================

\section*{Intentional notation overloads}

Only symbol-level overloads are recorded here.  The ambient mathematical type or explicit
adornment determines the intended sense.

\begin{longtable}{
@{}
>{\raggedright\arraybackslash}p{0.22\textwidth}
>{\raggedright\arraybackslash}p{0.50\textwidth}
>{\raggedright\arraybackslash}p{0.24\textwidth}
@{}
}

\toprule
\textbf{Notation} & \textbf{Permitted senses} & \textbf{Consistency rule}\\
\midrule
\endfirsthead
\toprule
\textbf{Notation} & \textbf{Permitted senses} & \textbf{Consistency rule}\\
\midrule
\endhead
\bottomrule
\endlastfoot

\CanonicalOverload
{\(f^*,\ f_*\)}
{For an algebra map \(A\to B\), these are extension and restriction of scalars on module categories.  For a geometric map \(f:Y\to X\), they are quasi-coherent pullback and its right adjoint; the pullback also acts on \(\operatorname{QCAlg}\).}
{Use the displayed domain and codomain.  On affines the geometric functors identify with extension and restriction of scalars.}

\CanonicalOverload
{\(\tau_{\geq n}^{\Sp},\ \tau_{\leq n}^{\Sp},\ \tau_{\geq n}^{A},\ \tau_{\leq n}^{A},\ \tau_{\geq n}^{X},\ \tau_{\leq n}^{X}\)}
{Postnikov truncations on spectra, on \(A\)-module categories, and on global quasi-coherent categories; \(\tau_{\geq0}^{\CAlg}\) is the separately adorned connective cover on commutative ring spectra.}
{Retain the ambient superscript when more than one Postnikov structure is in scope, and retain the \(\CAlg\) superscript for the algebra connective-cover functor.}

\CanonicalOverload
{\((-)_{\mathrm{cl}},\ X_{\mathrm{cl}}\)}
{Classical truncation is first defined on affines, then on spectral prestacks by \(\delta^*\), and finally restricted to spectral schemes.}
{These are compatible constructions; do not introduce separate notation for the three stages.}

\CanonicalOverload
{\(|X|\)}
{For an affine spectral scheme it is \(|\Spec\pi_0(A)|\); for a global spectral scheme it is \(|X_{\mathrm{cl}}|\).}
{The global definition extends the affine one.}

\CanonicalOverload
{\(-\otimes_A-\)}
{Relative tensor product in \(\Mod_A\); the same notation \(B\otimes_A C\) denotes the pushout of commutative \(A\)-algebras when the operands are algebra objects.}
{Use the algebra/module type of the operands; affine fibre products are represented by the algebra pushout.}

\CanonicalOverload
{\(\operatorname{Sym},\ \operatorname{Sym}_A,\ \operatorname{Sym}_{\mathcal O_X}\)}
{Free commutative algebra on spectra, on \(A\)-modules, and on quasi-coherent modules over \(X\), respectively.}
{Retain the base subscript once working relatively.}

\CanonicalOverload
{\(0,1\)}
{Ordinary numerical values throughout; in the affine-line construction they additionally denote the zero and unit sections \(X\rightrightarrows\mathbb A_X^1\).}
{The section sense is determined by the displayed source and target.}

\CanonicalOverload
{\(\mathbb L_X(\mathcal E),\ \mathbf V_X(\mathcal E)\)}
{\(\mathbb L_X(\mathcal E)\) is the relative linear space associated directly to \(\mathcal E\); \(\mathbf V_X(\mathcal E)=\mathbb L_X(\mathcal E^\vee)\) is the section total space for finite locally free \(\mathcal E\).}
{Preserve \(\mathbb L_X\) versus \(\mathbf V_X\): they represent dual functors of points and are not interchangeable.}

\CanonicalOverload
{\(\Sh,\ \Sh_\tau(\mathcal C)\)}
{In the earlier context-shape notation, bare \(\Sh\) denotes the fixed category of context shapes.  In this chapter the adorned expression \(\Sh_\tau(\mathcal C)\) denotes the \(\infty\)-category of sheaves for a topology \(\tau\).}
{Never use bare \(\Sh\) as an abbreviation for a sheaf category in the spectral chapter; retain both the topology subscript and the argument.}

\CanonicalOverload
{\(\Gamma,\ \Gamma(X,\mathcal F)\)}
{Earlier shaped syntax uses \(\Gamma\) as a context metavariable.  Here \(\Gamma(X,\mathcal F)\) is the persistent global-sections spectrum.}
{The parenthesized two-argument form is the global-sections notation; bare \(\Gamma\) remains available as a locally bound context or object variable according to scope.}

\CanonicalOverload
{\(\pi_0,\ \pi_n\)}
{Across the manuscript these symbols denote connected components or homotopy groups according to the ambient homotopical object; in this chapter they are used for spectra, ring spectra, and modules, with \(\pi_0\) also providing classical truncation of connective spectral algebras.}
{Use the ambient object type.  The ring structure on \(\pi_0(A)\) and module structures on \(\pi_n(M)\) are part of the spectral-algebraic use.}

\CanonicalOverload
{\(i_X,\ i_A,\ i\)}
{Here \(i_X:X_{\mathrm{cl}}^{\mathrm{disc}}\to X\) is the classical comparison morphism.  Earlier notation uses \(i_A\) for category of elements and bare \(i\) for the diagrammatic simplex inclusion.}
{Retain the \(X\)-subscript on the classical comparison morphism and the \(A\)-subscript on the category-of-elements functor.}

\CanonicalOverload
{\(\operatorname{Res}_A^B,\ \operatorname{Res}_\pi,\ \operatorname{Res}_{\pi,p}\)}
{In this chapter \(\operatorname{Res}_A^B\) is restriction of scalars along \(A\to B\).  Earlier shaped semantics uses \(\operatorname{Res}_\pi\) for ambient restriction and \(\operatorname{Res}_{\pi,p}\) for restricted realization.}
{Always retain the full subscript/superscript package; these are distinct constructions.}

\end{longtable}

% ============================================================================
% V. Intentional terminology overloads
% ============================================================================

\section*{Intentional terminology overloads}

The following technical words have established scoped senses at more than one level of the
spectral algebraic geometry construction.  They are not to be silently identified beyond
the compatibility statements proved in the chapter.

\begin{longtable}{
@{}
>{\raggedright\arraybackslash}p{0.22\textwidth}
>{\raggedright\arraybackslash}p{0.50\textwidth}
>{\raggedright\arraybackslash}p{0.24\textwidth}
@{}
}

\toprule
\textbf{Term} & \textbf{Permitted senses} & \textbf{Consistency rule}\\
\midrule
\endfirsthead
\toprule
\textbf{Term} & \textbf{Permitted senses} & \textbf{Consistency rule}\\
\midrule
\endhead
\bottomrule
\endlastfoot

\CanonicalTermOverload
{flat}
{Flat module spectrum; flat morphism of connective commutative ring spectra; flat morphism of spectral schemes; or flat quasi-coherent module.}
{Use the full noun phrase when more than one level is in scope.  Every sense is controlled by \(t\)-exact tensor/base change and is compatible with affine localization.}

\CanonicalTermOverload
{exact}
{Exact as a functor of stable \(\infty\)-categories versus left/right/fully \(t\)-exact for a Postnikov \(t\)-structure.}
{Do not use stable exactness as a synonym for geometric flatness; the chapter explicitly separates them.}

\CanonicalTermOverload
{connective}
{Connective spectrum, commutative ring spectrum, module, quasi-coherent module, or quasi-coherent algebra object.}
{Connectivity is always with respect to the relevant Postnikov \(t\)-structure and is detected on underlying spectra or affine restrictions as specified.}

\CanonicalTermOverload
{discrete}
{Discrete spectrum, discrete commutative ring spectrum, discrete affine spectral scheme, or discrete enhancement of an ordinary scheme.}
{All uses are Eilenberg--Mac Lane / degree-zero notions, but preserve the ambient object type.}

\CanonicalTermOverload
{open immersion}
{Affine spectral open immersion or representable open immersion of spectral Zariski stacks/schemes.}
{Global open immersions are defined by affine base change to the affine notion.}

\CanonicalTermOverload
{closed immersion}
{Affine closed immersion, detected by surjectivity on \(\pi_0\), or global closed immersion tested after affine base change.}
{Global closedness is affine-local and classically detected.}

\CanonicalTermOverload
{affine}
{Affine spectral scheme; affine open immersion; or affine morphism of spectral schemes.}
{State the noun being modified; an affine morphism is characterized by affine pullback over every affine open of the target.}

\CanonicalTermOverload
{finite presentation / finite cell}
{``Locally of finite presentation'' is the compactness condition; ``finite cell'' is the explicit finite-colimit closure of finitely generated free commutative algebras.}
{Do not identify the terms.  A locally finitely presented connective algebra is a retract of a finite cell algebra.}

\CanonicalTermOverload
{classical truncation}
{Affine truncation, prestack restriction to discrete affines, and the resulting functor on spectral schemes.}
{The constructions are compatible and use the same notation \((-)_{\mathrm{cl}}\).}

\CanonicalTermOverload
{vector bundle / locally free}
{Locally free modules on affines and locally free quasi-coherent modules globally; a vector bundle of rank \(n\) is the global fixed-rank notion.}
{Finite projective is the affine intrinsic characterization; locally constant rank need not be globally constant.}

\end{longtable}

\medskip
\noindent
\emph{Projection rule.}
The publication-facing List of Notation is a lean navigation projection drawn only from
entries marked ``yes'' above.  Every canonical head notation item displayed there is backed
by a public entry in this register.  The public display may append registered type, signature,
or defining-equation information, and tightly coupled aliases or adjoint pairs may be merged
into one row, but no such presentation change alters a symbol's registered mathematical type
or technical meaning.  New persistent notation and terminology are registered here first;
the public list is then updated from this register.

\endgroup
